% Content-only chapter block for TFM.tex
% Included from: chapters/01-06-introduction-and-research-design/chapter.tex

% !TEX program = xelatex
% !TEX root = ../../TFM.tex
\documentclass[../../TFM.tex]{subfiles}

\begin{document}

	% =====================================================
	% INTRODUCCIÓN
	% =====================================================
	\chapter{Introducción}
	\label{chap:introduccion}
	\thesisstartchap{introduccion}

	Los recursos naturales han ocupado históricamente un lugar central en las trayectorias de desarrollo de numerosos países. Las primeras interpretaciones estructuralistas señalaron que la división internacional del trabajo tendía a especializar a las economías periféricas en la exportación de bienes primarios, condicionando así sus patrones de crecimiento y sus procesos de industrialización. En este sentido, \textcite{Prebisch1962} sostuvo que la estructura centro-periferia de la economía mundial limitaba las posibilidades de desarrollo de los países exportadores de materias primas, mientras que \textcite{Singer1975} destacó la tendencia al deterioro de los términos de intercambio de los productos primarios frente a las manufacturas.

	Sin embargo, la evidencia histórica también muestra que los recursos naturales han servido de base para experiencias exitosas de desarrollo. Estudios sobre desarrollo basado en recursos sostienen que su abundancia puede favorecer la acumulación de capacidades productivas, el aprendizaje tecnológico y el crecimiento de largo plazo cuando existen políticas e instituciones adecuadas \parencite{Wright2004,Ville2013}. En consecuencia, la relación entre recursos naturales y desarrollo sigue siendo objeto de debate, particularmente en torno a las condiciones bajo las cuales estos recursos actúan como restricción estructural o como plataforma para la transformación productiva.

	A partir de la década de 1990, una parte importante de la literatura comenzó a cuestionar la idea de que los recursos naturales constituyan necesariamente una ventaja para el desarrollo. En este marco, la denominada hipótesis de la ``maldición de los recursos''\footnote{El término ``maldición de los recursos'' se utiliza aquí como una categoría analítica de la literatura, no como una relación determinista entre recursos naturales y bajo desempeño económico. La tesis parte precisamente de que los efectos de los recursos dependen de condiciones productivas, institucionales y macroeconómicas.} sostiene que las economías con alta dependencia de exportaciones primarias o con fuerte peso macroeconómico de las rentas extractivas tienden a experimentar menores tasas de crecimiento y mayores dificultades para sostener procesos de transformación productiva.

	Uno de los trabajos más influyentes en este debate es el de \textcite{Sachs2001}, quienes encuentran una relación negativa entre la especialización en recursos y el crecimiento de largo plazo. Esta perspectiva retoma contribuciones tempranas como la de \textcite{Auty1993}, quien argumenta que las economías intensivas en recursos enfrentan incentivos estructurales que pueden obstaculizar la industrialización. Desde una perspectiva de economía política, \textcite{Ross1999} subraya que las rentas provenientes de los recursos pueden afectar la calidad institucional y la rendición de cuentas, generando dinámicas que limitan el desarrollo sostenido.

	En conjunto, esta literatura puso de relieve que los efectos económicos de los recursos naturales no son unívocos y que resulta necesario distinguir entre abundancia, dependencia externa y peso económico de las rentas\footnote{La abundancia refiere a la dotación física o geológica de recursos; la dependencia externa al peso de los recursos en la estructura exportadora; y las rentas extractivas al ingreso económico generado por la explotación de esos recursos en relación con el tamaño de la economía. Esta distinción se desarrolla formalmente en el capítulo metodológico.} para comprender sus implicaciones sobre el crecimiento y la transformación estructural \parencite{VanDerPloeg2011}.

	Diversos estudios han buscado identificar los mecanismos económicos e institucionales que podrían explicar estos resultados. Desde una perspectiva macroeconómica, uno de los canales más discutidos es la denominada enfermedad holandesa (\textit{Dutch disease}), mediante la cual un auge en el sector de recursos naturales puede provocar una apreciación del tipo de cambio real y un desplazamiento de factores productivos hacia sectores no transables, debilitando así el desarrollo del sector manufacturero y de otras actividades transables \parencite{Corden1982}. Asimismo, la volatilidad de los precios internacionales de los commodities puede generar ciclos de auge y caída que afectan la estabilidad macroeconómica y dificultan la formulación de políticas contracíclicas en economías dependientes de recursos \parencite{Anne2021}.

	Más allá de estos efectos macroeconómicos, la literatura también ha destacado dinámicas de economía política vinculadas a la competencia por las rentas derivadas de los recursos. En este sentido, \textcite{Tornell1999} introducen el concepto de \textit{voracity effect} para describir cómo la competencia entre grupos de interés por apropiarse de dichas rentas puede conducir a una expansión excesiva del gasto y a una asignación ineficiente de los recursos.

	La importancia de las instituciones en este proceso es subrayada por \textcite{Mehlum2006}, quienes argumentan que los efectos de los recursos dependen crucialmente de si las instituciones incentivan actividades productivas o, por el contrario, fomentan comportamientos de búsqueda de rentas. En línea con esta perspectiva, evidencia empírica sugiere que la disponibilidad de recursos puede intensificar problemas de corrupción y debilitar la calidad institucional, afectando negativamente el desempeño económico de largo plazo \parencite{Bhattacharyya2010}. En conjunto, estos aportes sugieren que los efectos de los recursos sobre el desarrollo y la transformación estructural dependen de la interacción entre restricciones macroeconómicas y la calidad del marco institucional en el que se gestionan las rentas extractivas.

	En este contexto, una parte importante de la literatura ha centrado su atención en la relación entre la dependencia de recursos naturales y las dinámicas de cambio estructural. Diversos estudios señalan que las economías intensivas en recursos pueden enfrentar dificultades para desarrollar sectores manufactureros y otras actividades transables, particularmente en un contexto de desindustrialización temprana que afecta a muchas economías en desarrollo \textcite{Rodrik2016}.

	No obstante, la evidencia empírica también sugiere que los resultados asociados a la especialización en recursos dependen en gran medida de la capacidad de las economías para acumular capacidades, densificar encadenamientos productivos y diversificar su estructura productiva. En esta línea, \textcite{Agosin2012} analizan los determinantes de la diversificación exportadora y encuentran que la acumulación de capacidades productivas y el desarrollo institucional desempeñan un papel central en la evolución de la estructura productiva. Desde una perspectiva de comercio internacional, \textcite{Harding2016} muestran que las exportaciones de recursos naturales pueden influir significativamente en la composición del comercio no basado en recursos, generando tanto oportunidades como restricciones para el desarrollo de otras actividades transables.

	De manera complementaria, trabajos sobre industrialización basada en recursos destacan que, bajo ciertas condiciones, los auges de commodities pueden convertirse en plataformas para el desarrollo de nuevas actividades productivas a través de encadenamientos, aprendizaje tecnológico y expansión industrial \parencite{Morris2013}. En este sentido, el desafío para las economías dependientes de recursos no radica únicamente en la presencia de recursos abundantes, sino en la capacidad de transformar las rentas extractivas en procesos sostenidos de transformación estructural y acumulación de capacidades.

	En este marco, una literatura creciente ha puesto el énfasis en el papel de las capacidades productivas y de la estructura del conocimiento en los procesos de desarrollo. Desde esta perspectiva, el crecimiento sostenido no depende únicamente de la disponibilidad de recursos, sino de la capacidad de las economías para acumular capacidades y avanzar hacia actividades más sofisticadas. En este sentido, \textcite{Hausmann2014} sostienen que el desarrollo económico está estrechamente vinculado a la complejidad de la estructura productiva, entendida como el conjunto de conocimientos y capacidades incorporados en las actividades económicas de un país.

	De manera complementaria, el concepto de \textit{product space} desarrollado por \textcite{Hausmann2007} destaca que la posibilidad de expandir la estructura productiva depende de la proximidad entre las capacidades existentes y aquellas requeridas por actividades más sofisticadas. Estas ideas retoman, en parte, intuiciones tempranas sobre la importancia de los encadenamientos productivos y los procesos de aprendizaje en la transformación estructural de las economías \parencite{Hirschman1958}.

	Más recientemente, estudios sobre economías dependientes de recursos han subrayado que la capacidad de convertir las rentas provenientes de los recursos naturales en procesos sostenidos de acumulación de capacidades y diversificación productiva depende tanto del desarrollo previo de dichas capacidades como de las estrategias institucionales y de política que faciliten la transición hacia actividades más complejas \parencite{Lebdioui2019}. En este contexto, comprender bajo qué condiciones las economías intensivas en recursos logran traducir sus rentas extractivas en transformación estructural constituye una cuestión central para el análisis del desarrollo económico contemporáneo.

	A partir de este marco conceptual, la presente investigación examina las
	asociaciones entre las rentas extractivas, distintos canales teóricos y la
	transformación estructural externa de economías dependientes de recursos
	naturales. El estudio distingue entre la dependencia externa de recursos
	---utilizada únicamente para delimitar la muestra a partir de la estructura
	exportadora--- y la intensidad macroeconómica de las rentas extractivas,
	incorporada como variable explicativa. La transformación estructural externa
	se aproxima principalmente mediante la complejidad económica y, de manera
	complementaria, mediante la diversificación de la canasta exportadora. El
	diseño permite estimar asociaciones condicionales dentro de los países, pero
	no identificar efectos causales ni observar directamente la conversión de las
	rentas en capacidades productivas.

	En lugar de evaluar si la abundancia de recursos constituye en sí misma una ventaja o una desventaja para el crecimiento económico, el análisis compara cómo los canales productivos, institucionales y macroeconómicos se relacionan con $ECI$ y $DIVX$ dentro del universo seleccionado. Para ello, el estudio adopta un enfoque cuantitativo longitudinal. El período 1990--1995 se utiliza exclusivamente para seleccionar 55 economías con $DRES\geq20\%$. La estimación se realiza para 1996--2021 sobre una muestra común desbalanceada de 1.044 observaciones, 49 países y 23 años efectivos. Los insumos originales cubren horizontes más amplios según la fuente, pero no definen una muestra econométrica adicional.

	En este contexto, la investigación busca contribuir a la literatura sobre desarrollo económico y economías extractivas mediante un análisis comparado que vincula dependencia externa de recursos, rentas extractivas, transformación estructural y acumulación de capacidades productivas. Tras presentar la pregunta, la justificación, el planteamiento del problema, los objetivos y las hipótesis, la \autoref{sec:contribucion} precisa los aportes esperados de la investigación. El resto del documento se organiza de la siguiente manera: el \autoref{chap:marco-teorico} desarrolla el marco teórico y conceptual del estudio; el \autoref{chap:literatura-empirica} revisa la literatura empírica relevante; el \autoref{chap:metodologia} presenta la estrategia metodológica y las técnicas empíricas utilizadas; el \autoref{chap:resultados} expone los resultados del análisis econométrico; y finalmente el \autoref{chap:conclusiones} discute las implicaciones de los hallazgos y las principales conclusiones de la investigación.

\thesisendchap{introduccion}
	% =====================================================
	% PREGUNTA
	% =====================================================
	\chapter{Pregunta}
	\label{chap:pregunta}
	\thesisstartchap{pregunta}

	La literatura sobre economías dependientes de recursos naturales ha destacado la existencia de trayectorias de desarrollo divergentes. Mientras algunas economías logran avanzar hacia estructuras exportadoras más complejas y diversificadas, otras permanecen atrapadas en patrones de especialización extractiva.

	En este contexto, la pregunta general que guía esta investigación es:

	\begin{quote}
	¿Cómo se asocian las rentas extractivas y los canales productivos,
	institucionales y macroeconómicos con la transformación estructural externa de
	las economías dependientes de recursos naturales?\footnote{La transformación
	estructural externa se observa principalmente mediante el $ECI$, como
	aproximación a la sofisticación y a las capacidades productivas reveladas, y
	complementariamente mediante $DIVX=1-HHI$, como medida de diversificación de
	la canasta exportadora.}
	\end{quote}

	Con el fin de abordar empíricamente esta cuestión, la tesis se propone responder la siguiente pregunta específica:

	\begin{quote}
	¿Qué asociaciones condicionales se observan durante 1996--2021 entre las
	rentas extractivas, la calidad institucional, la estructura exportadora, las
	condiciones macroeconómicas, las capacidades productivas y los resultados de
	complejidad económica y diversificación exportadora, una vez controlados los
	efectos fijos de país y año?
	\end{quote}

\thesisendchap{pregunta}
	% =====================================================
	% JUSTIFICACIÓN
	% =====================================================
	\chapter{Justificación}
	\label{chap:justificacion}
	\thesisstartchap{justificacion}

	La elección del presente tema se justifica por su relevancia teórica y empírica para la comprensión de las trayectorias de desarrollo de economías dependientes de recursos naturales. La persistencia de estructuras exportadoras altamente concentradas en actividades extractivas, junto con la marcada heterogeneidad observada en los procesos de crecimiento y transformación estructural entre países, constituye un problema central del análisis económico contemporáneo. Estas dinámicas se asocian frecuentemente con elevados niveles de volatilidad macroeconómica, bajos niveles de complejidad económica, escasa diversificación exportadora y dificultades para la acumulación sostenida de capacidades productivas en el largo plazo \parencite{Auty1993,Sachs1995}.

	Desde una perspectiva académica, la investigación se inserta en un debate amplio y aún abierto sobre los efectos económicos de la dependencia de recursos naturales y de las rentas extractivas. Si bien la literatura ha avanzado de manera significativa en la identificación de mecanismos vinculados a la denominada ``maldición de los recursos'', los resultados empíricos muestran una notable diversidad de trayectorias \parencite{Brunnschweiler2006,Brunnschweiler2008,Lebdioui2019}. Mientras algunas economías han logrado traducir las rentas extractivas en mejoras de complejidad económica y diversificación exportadora, otras permanecen atrapadas en trayectorias de especialización primaria \parencite{Lebdioui2019,Ville2013}. Esta coexistencia de resultados divergentes pone de manifiesto la necesidad de profundizar el análisis de las condiciones productivas, institucionales y macroeconómicas que median entre la dependencia extractiva, la gestión de las rentas y los resultados de desarrollo observados.

	\textcite{Ajide2022} señalan que, aunque la literatura sobre la maldición de los recursos naturales se ha centrado predominantemente en sus efectos sobre el crecimiento económico, su impacto sobre dimensiones estructurales del desarrollo, como la complejidad económica, ha sido ampliamente ignorado. En este sentido, la complejidad económica ofrece una aproximación relevante a la transformación estructural externa, al captar la diversificación y el contenido tecnológico de la canasta exportadora. Para los fines de este trabajo, ese énfasis se complementa con una medida directa de diversificación exportadora, definida como el complemento de la concentración de Herfindahl--Hirschman. De este modo, la investigación busca cubrir un vacío de la literatura empírica mediante una aproximación centrada en la sofisticación y diversificación de la inserción productiva externa.

	En este contexto, resulta pertinente adoptar un enfoque integrador que permita examinar de manera conjunta los mecanismos estructurales que condicionan dichas trayectorias. El análisis comparativo entre países, basado en indicadores de dependencia extractiva, rentas extractivas, complejidad económica, diversificación exportadora, capacidades humanas, innovación y desempeño macroeconómico, ofrece una vía adecuada para identificar regularidades y contrastar hipótesis sobre factores asociados con el cambio estructural. De este modo, la investigación busca contribuir a una comprensión más precisa de los canales teóricos vinculados con las diferencias observadas en la complejidad y la diversificación exportadoras de las economías basadas en recursos naturales, evitando enfoques deterministas y simplificaciones normativas.

\thesisendchap{justificacion}

	% PLANTEAMIENTO DEL TEMA / PROBLEMA
	% =====================================================
	\chapter{Planteamiento del tema / problema}
	\label{chap:planteamiento}
	\thesisstartchap{planteamiento}

	La abundancia de recursos naturales ha sido tradicionalmente concebida como una fuente potencial de crecimiento y desarrollo económico. Sin embargo, la evidencia empírica muestra que numerosas economías altamente dependientes de actividades extractivas presentan trayectorias caracterizadas por baja complejidad económica, elevada volatilidad macroeconómica y persistencia de estructuras productivas poco diversificadas y concentradas en sectores primarios \parencite{Tornell1999,Arezki2012}. Al mismo tiempo, existen casos en los que las rentas provenientes de recursos naturales han sido transformadas en procesos sostenidos de acumulación de capacidades, transformación productiva y fortalecimiento institucional \parencite{Ville2013}.

	Esta heterogeneidad plantea un problema central: la dependencia de recursos naturales no conduce de manera automática ni a la denominada ``maldición'' ni al desarrollo exitoso, sino que se relaciona con resultados heterogéneos según la estructura productiva, el marco institucional y las condiciones macroeconómicas \parencite{Lebdioui2019}. En particular, persiste una brecha analítica sobre la forma en que las rentas extractivas y esos canales se asocian con la sofisticación y la diversificación exportadoras dentro de un conjunto comparable de economías dependientes.

	El problema de investigación consiste en evaluar comparativamente esas asociaciones durante 1996--2021. La dependencia exportadora de 1990--1995 delimita el universo de análisis; dentro de este, las rentas extractivas y los canales institucional, estructural, macroeconómico, de capacidades productivas, fiscal y financiero se incorporan como variables explicativas de $ECI$ y $DIVX$.

	Este planteamiento permite contrastar hipótesis de asociación mediante modelos con efectos fijos de país y año estimados sobre una muestra común. La investigación busca identificar regularidades dentro de los países y comparar su estabilidad entre $ECI$ y $DIVX$, sin interpretar los coeficientes como efectos causales ni como mediciones directas de mecanismos no observados.

\thesisendchap{planteamiento}
	% =====================================================
	% OBJETIVOS
	% =====================================================

	\chapter{Objetivos}
	\label{chap:objetivos}
	\thesisstartchap{objetivos}

	\section{Objetivo general}
	\label{sec:objetivo-general}



	Analizar comparativamente cómo las rentas extractivas y los canales
	productivos, institucionales y macroeconómicos se asocian con la complejidad
	económica y la diversificación exportadora de las economías dependientes de
	recursos naturales durante 1996--2021, mediante modelos de panel con efectos
	fijos de país y año.



	\section{Objetivos específicos}

	\label{sec:objetivos-especificos}



	\begin{enumerate}



		\item Delimitar una muestra de economías dependientes de recursos naturales
		a partir de $DRES$ durante 1990--1995 y distinguir ese criterio de selección
		de las rentas extractivas observadas durante el período econométrico.



		\item Estimar la asociación de $RENTS$ con $ECI$ y $DIVX$ y evaluar si la
		interacción $RENTS\times INST$ modifica esa asociación dentro del soporte
		observado de la calidad institucional.



		\item Examinar cómo la concentración y la composición exportadoras, las
		rentas extractivas específicas por habitante y la desagregación entre
		hidrocarburos y minería se relacionan con los dos resultados de
		transformación estructural externa.



		\item Evaluar las asociaciones de la volatilidad de los términos de
		intercambio y del tipo de cambio real con $ECI$ y $DIVX$, junto con los
		controles macroeconómicos incluidos en la especificación completa.



		\item Evaluar en qué medida el capital humano, la innovación y la
		conectividad digital se asocian con $ECI$ y $DIVX$, controlando por nivel de
		desarrollo, consumo público y profundidad financiera.



		\item Evaluar la selección del estimador de panel y examinar la estabilidad
		de los resultados mediante pruebas de estacionariedad, primeras diferencias,
		sensibilidad dinámica, inferencia bootstrap y Driscoll--Kraay, exclusión
		iterativa de países y observaciones influyentes, tendencias lineales por país,
		sensibilidad temporal de 2014 y desagregación de las rentas.



	\end{enumerate}


	\thesisendchap{objetivos}

	% =====================================================
	% HIPÓTESIS
	% =====================================================
	\chapter{Hipótesis}

	\label{chap:hipotesis}

	\thesisstartchap{hipotesis}



	La presente investigación se guía por hipótesis de asociación derivadas de los
	mecanismos productivos, institucionales y macroeconómicos discutidos en el
	\autoref{chap:marco-teorico}. Las hipótesis se refieren a signos y patrones
	esperados en $ECI$ y $DIVX$; no convierten el diseño observacional en una
	estrategia de identificación causal.






	\section{Hipótesis general}

	\label{sec:hipotesis-general}



	Dentro de las economías dependientes de recursos naturales, una mayor
	intensidad de las rentas extractivas se asocia con menor complejidad económica
	y menor diversificación exportadora, aunque la magnitud de esta relación puede
	variar según la calidad institucional, la estructura exportadora, las
	condiciones macroeconómicas y las capacidades productivas observadas.


	\section{Hipótesis específicas}
	\label{sec:hipotesis-especificas}

	\begin{enumerate}
		\item \label{hyp:h1} Mayores rentas extractivas y una estructura exportadora
		más concentrada o especializada en productos primarios se asocian con menor
		$ECI$ y menor $DIVX$.\footnote{La expresión ``se asocia con'' refiere a una
		relación condicional estimada con datos de panel y no a un efecto causal.}

		\item \label{hyp:h2} Una mayor calidad institucional atenúa la asociación
		negativa de $RENTS$ con $ECI$ y $DIVX$, por lo que se espera un efecto
		marginal menos adverso de las rentas cuando $INST$ es más alto.

		\item \label{hyp:h3} Una mayor volatilidad de los términos de intercambio de
		commodities y una apreciación del tipo de cambio real se asocian con menores
		niveles de $ECI$ y $DIVX$.

		Esta hipótesis se contrasta mediante $VOL$ y $RER$. Ambos indicadores aproximan
		condiciones macroeconómicas observables; no identifican directamente decisiones
		de ahorro, inversión o gestión fiscal procíclica.

		\item \label{hyp:h4} Mayores niveles de capital humano e innovación se
		asocian positivamente con $ECI$ y $DIVX$.

		\item \label{hyp:h5} La asociación de las rentas extractivas con $ECI$ y
		$DIVX$ difiere entre el componente de petróleo y gas y el componente de
		minería y carbón.
	\end{enumerate}


	\section{Contribución de la investigación}
	\label{sec:contribucion}

	La contribución de esta investigación no consiste en proponer una teoría completamente nueva sobre la relación entre recursos naturales y desarrollo, ni en afirmar que cada variable utilizada sea inédita en la literatura. Existe ya una línea empírica reciente que examina vínculos muy próximos entre recursos naturales, rentas extractivas y complejidad económica, tanto en muestras regionales como en paneles amplios de países \parencite{Ajide2022,Avom2022,Canh2020,Owjimehr2024,Zhang2023}. En diálogo con esos trabajos, el aporte de esta tesis radica en reorganizar empíricamente el problema de la transformación estructural en economías dependientes de recursos naturales no renovables, distinguiendo entre el criterio que define el universo de análisis, las variables explicativas y los resultados analizados. En ese sentido, el trabajo no se presenta como una réplica ni como una sustitución de dicha literatura, sino como una aplicación comparada que evalúa asociaciones entre canales y dos dimensiones distintas de la inserción exportadora.

	El primer aporte consiste en separar la dependencia externa de recursos naturales de la intensidad macroeconómica de las rentas extractivas. Varios estudios analizan la relación entre rentas de recursos naturales y complejidad económica en muestras amplias de países, países en desarrollo o economías ricas en recursos, y algunos distinguen incluso entre petróleo, minería y gas \parencite{Canh2020,Hoang2023,Li2024,Ul-Durar2023,Valentine2024,Zhang2023}. Esta tesis toma ese punto de partida, pero utiliza la dependencia externa, medida por la participación de las exportaciones de recursos naturales no renovables del subsuelo ---hidrocarburos, carbón, minerales y metales--- en la canasta exportadora, como criterio para delimitar la muestra. En cambio, las rentas de recursos naturales como porcentaje del PIB se incorporan como una variable explicativa dentro del modelo econométrico. Esta distinción evita confundir la dependencia de recursos como condición estructural de partida con las rentas extractivas cuya asociación condicional se estima durante el período econométrico.

	El segundo aporte se encuentra en las variables explicadas. La literatura más cercana se concentra principalmente en el Índice de Complejidad Económica (ECI) como resultado para estudiar si las rentas, la abundancia o la dependencia de recursos naturales limitan la sofisticación productiva de las economías \parencite{Ajide2022,Canh2020,Hoang2023,Owjimehr2024,Valentine2024,Zhang2023}. En este marco, la tesis aproxima la transformación estructural desde la inserción externa. La variable dependiente principal es el ECI, entendido como una medida de sofisticación exportadora y de capacidades productivas reveladas. De manera complementaria, se incorpora la diversificación exportadora, definida como $DIVX=1-HHI$, donde $HHI$ corresponde al índice de concentración de Herfindahl--Hirschman. Esta segunda variable no reemplaza al ECI, sino que permite observar si los mecanismos asociados a la complejidad económica también se relacionan con una canasta exportadora menos concentrada. Así, la tesis distingue entre sofisticación y diversificación, dos dimensiones próximas pero no equivalentes de la transformación estructural externa.

	El tercer aporte consiste en articular en un mismo marco empírico distintos canales explicativos que la literatura cercana suele tratar de manera separada. Algunos estudios enfatizan el papel de las instituciones, la estabilidad política o la calidad institucional en economías ricas en recursos \parencite{Ajide2022,Avom2022,Olaniyi-Odhiambo2025-3,Olaniyi-Odhiambo2025,Owjimehr2024}; otros destacan el desarrollo financiero, la innovación, el capital humano o el nivel de ingreso como condiciones que pueden modificar la relación entre recursos y complejidad \parencite{Alvarado2023,Li2024,Valentine2024}. A partir de estos antecedentes, el modelo incorpora un canal institucional, asociado a la calidad institucional y a su interacción con las rentas extractivas; rentas específicas de petróleo, gas y carbón por habitante; un canal estructural, relacionado con concentración y especialización exportadora; un canal macroeconómico, asociado a volatilidad y tipo de cambio real; un canal de capacidades productivas, aproximado mediante capital humano, innovación y conectividad; y controles fiscales y financieros. $M_1$ y $M_2$ organizan estos canales en especificaciones temáticas paralelas, mientras $M_3$ los integra en el modelo completo. El aporte no reside en afirmar que cada canal sea novedoso, sino en compararlos sobre la misma muestra y frente a dos resultados no equivalentes.

	Finalmente, la tesis aporta una estrategia empírica comparada que busca diferenciarse de los estudios existentes no por el simple uso de datos de panel, sino por la forma en que define y organiza el universo de análisis. En efecto, ya existen aplicaciones de panel sobre recursos naturales y complejidad económica, incluyendo estudios con muestras globales, países en desarrollo, economías ricas en recursos y casos africanos \parencite{Ajide2022,Avom2022,Owjimehr2024,Valentine2024,Zhang2023}, así como análisis de causalidad o no linealidad entre complejidad y rentas de recursos naturales \parencite{Olaniyi-Odhiambo2025-2,Ul-Durar2023}. Frente a esos antecedentes, el diseño selecciona 55 países cuya dependencia externa de recursos alcanza al menos 20\% durante 1990--1995 y estima, para 1996--2021, un modelo principal con $ECI$ y una especificación complementaria con $DIVX=1-HHI$. La muestra econométrica común contiene 1.044 observaciones de 49 países y 23 años efectivos. En conjunto, la contribución es ofrecer evidencia comparada sobre asociaciones condicionales, sin atribuir causalidad al diseño ni presentar como inéditos elementos que la literatura ya ha trabajado.

	\thesisendchap{hipotesis}

\end{document}


